\setchapterpreamble[u]{\margintoc}
\chapter{Estimation}
\labch{estimation}

\sources{Bertsekas and Tsitsiklis, \emph{Introduction to Probability}, 2nd ed.\
\cite{bertsekas2008probability}, Chapters~8--9; MIT 6.041 \cite{ocw6041},
Lectures~21--25; MIT 18.05 \cite{ocw1805}.}

Estimation is where probability meets the fitted model. This chapter shows that
the two losses used almost universally in supervised learning are not arbitrary
choices but maximum likelihood under two different noise assumptions --- and
that weight decay is a prior.

\section{The problem}
\labsec{estimation-problem}

Data \(x_1,\ldots,x_n\) are assumed drawn from a family of distributions indexed
by an unknown parameter \(\theta\). An \emph{estimator} is any function
\(\hat\theta(x_1,\ldots,x_n)\) of the data. It is a random variable, so it has a
bias \(\E[\hat\theta]-\theta\) and a variance, and its quality is usually
measured by mean squared error, which decomposes as
\[
	\E\bigl[(\hat\theta-\theta)^{2}\bigr]
	=\bigl(\E[\hat\theta]-\theta\bigr)^{2}+\Var(\hat\theta).
\]

\marginnote{Bias and variance trade against each other. The regularisation of
\refsec{regularisation} buys a large reduction in the second term at the cost of
a small increase in the first --- which is why it works.}

\section{Maximum likelihood}
\labsec{mle}

\begin{definition}
The \emph{likelihood} is the joint density of the observed data read as a
function of \(\theta\),
\[
	L(\theta)=\prod_{i=1}^{n}f(x_i\mid\theta),
\]
and the maximum likelihood estimate is the \(\theta\) maximising it, equivalently
maximising \(\log L(\theta)=\sum_i\log f(x_i\mid\theta)\).
\end{definition}

The logarithm is taken because it converts the product into a sum, which
differentiates cleanly and does not underflow.

\begin{example}
For \(X_i\sim\)~Bernoulli\((p)\),
\(\log L(p)=\sum_i\bigl[x_i\log p+(1-x_i)\log(1-p)\bigr]\). Differentiating and
setting to zero gives \(\hat p=\bar x\), the observed frequency.
\end{example}

\begin{example}
For \(X_i\sim\)~Gaussian\((\mu,\sigma^{2})\) with \(\sigma\) known,
\[
	\log L(\mu)=-\frac{1}{2\sigma^{2}}\sum_i(x_i-\mu)^{2}+\text{const},
\]
so \(\hat\mu=\bar x\). Estimating \(\sigma^{2}\) as well gives
\(\hat\sigma^{2}=\frac1n\sum_i(x_i-\bar x)^{2}\), which is biased low; dividing by
\(n-1\) instead removes the bias.
\end{example}

\section{Two losses, one principle}
\labsec{losses-from-likelihood}

Now let the parameter be the weights of a model \(f_{\vect{w}}\) and condition on
inputs.

\paragraph{Gaussian noise gives squared error.}
Assume \(y_i=f_{\vect{w}}(\vect{x}_i)+\varepsilon_i\) with
\(\varepsilon_i\sim\)~Gaussian\((0,\sigma^{2})\), independent. Then
\[
	\log L(\vect{w})
	=-\frac{1}{2\sigma^{2}}\sum_{i=1}^{n}\bigl(y_i-f_{\vect{w}}(\vect{x}_i)\bigr)^{2}
	+\text{const},
\]
so maximising the likelihood is minimising the sum of squared errors. With a
linear model this is precisely the least-squares problem of
\refch{orthogonality}.

\paragraph{Bernoulli labels give cross-entropy.}
Assume \(y_i\in\{0,1\}\) with
\(\Prob(y_i=1)=\hat y_i=f_{\vect{w}}(\vect{x}_i)\). Then
\[
	-\log L(\vect{w})
	=-\sum_{i=1}^{n}\bigl[y_i\log\hat y_i+(1-y_i)\log(1-\hat y_i)\bigr],
\]
which is the cross-entropy loss.

\begin{remark}
A question left open in \refsec{worked-gradient} is settled here. Squared
error applied to a logistic output carries a factor \(\hat y(1-\hat y)\) that vanishes when the unit
saturates; cross-entropy is the loss that the Bernoulli assumption actually
implies, and its gradient with respect to the pre-activation reduces to
\(\hat y-y\), with no vanishing factor. The improvement is a consequence of
matching the loss to the noise model rather than a numerical trick.
\end{remark}

\section{Bayesian estimation}
\labsec{bayesian}

Maximum likelihood treats \(\theta\) as fixed and unknown. The Bayesian
alternative gives it a distribution: a \emph{prior} \(p(\theta)\) before the data
and a \emph{posterior} after, related by \refeq{bayes},
\[
	p(\theta\mid \text{data})\ \propto\ L(\theta)\,p(\theta).
\]
Reporting the maximiser of the posterior is \emph{maximum a posteriori}
estimation.

\begin{proposition}
For a linear model with Gaussian noise of variance \(\sigma^{2}\) and a Gaussian
prior \(\vect{w}\sim\)~Gaussian\((\vect{0},\tau^{2}I)\), the maximum a posteriori
estimate minimises
\[
	\frac{1}{2\sigma^{2}}\lVert X\vect{w}-\vect{y}\rVert^{2}
	+\frac{1}{2\tau^{2}}\lVert\vect{w}\rVert^{2},
\]
which is ridge regression with \(\delta^{2}=\sigma^{2}/\tau^{2}\).
\end{proposition}

\begin{proof}
Take logarithms in the display above. The log-likelihood contributes the first
term up to a constant, and the log-prior of a Gaussian contributes the second.
Maximising the sum is minimising its negative.
\end{proof}

\marginnote{So the \(\delta^{2}I\) added in \refsec{regularisation} for numerical
reasons, the weight decay added in training for empirical reasons, and a Gaussian
prior on the weights are three descriptions of one modification.}

\section{What this part establishes}
\labsec{part3-summary}

Three conclusions carry into Part~IV. Averages over a sample estimate
expectations, with error shrinking like \(1/\sqrt n\). Squared error and
cross-entropy are maximum likelihood under Gaussian and Bernoulli assumptions
respectively. And regularisation is a prior, which means the choice of penalty is
a statement of belief about the parameters rather than a free knob.
